\section{Signs and removal of the disk obstruction}
\label{obs:section}

The finite integrations of Theorem~\ref{virt:finite-integration} provide
geometric correspondences.  We now specify their signs and use the
correspondence with one interior mark to remove curvature.  The topological
input from Proposition~\ref{top:brane} is precisely that
\begin{equation}
 \label{obs:topological-input}
 i_{L*}:H_2(L;\Q)\longrightarrow H_2(T;\Q),\qquad
 i_{L*}:H_6(L;\Q)\longrightarrow H_6(T;\Q)
 \quad\text{are injective}.
\end{equation}
Together with the grading modulo four, these two injections suffice; no
injectivity in the other homological degrees is required.

Throughout this section all sign exponents are read modulo two.  For an
operation whose boundary inputs are written counterclockwise after its
output, the suspended degree of an input $x$ is $|x|-1$.  We write
$A^{\bar r}(L)$ for forms whose ordinary degrees are congruent to $r$
modulo four.  Formal Taylor variables have degree zero.  The notation
$\equiv_N$ means equality after discarding terms of energy greater than
$N$, coefficient by coefficient in the Taylor variable.

\subsection{Cappings, degrees, and local gluing orientations}

At a transverse intersection choose capping operators joining the two
graded tangent planes, with matching lifts of their gradings.  Use the
reverse path for the other end of the same generator.  Choose spin frames
on the two branes at the intersection, extend them across the cap, and
use the reversed frames for the reciprocal cap.  On a brane boundary arc
use a spin trivialization matching its endpoint frames.  A basis of each
cap determinant line is fixed up to positive scaling.  These are the
ordinary spin cappings with background class zero; we use their
determinant gluing and excision constructions
\cite[Section~5]{akaho-joyce}.  The particular comparisons needed below are
spelled out, including their order of factors.

\begin{lemma}[Degrees]\label{obs:degree}
The operation $m_k$ has degree $2-k$ modulo four, and has that degree
integrally on a strictly directed graph word.  A pure disk of Maslov
index $\mu$ with $k$ boundary inputs has virtual dimension
\[
 8+\mu+k-2
 \quad\text{with a boundary output},\qquad
 8+\mu+k-1
 \quad\text{with an interior output}.
\]
In particular the curvature of a cochain in $A^{\bar1}(L)$ has only
ordinary form degrees $2$ and $6$.
\end{lemma}

\begin{proof}
The cap attached to an input has a negative end and has index equal to
the generator degree.  To check this convention, split a transverse
plane pair into positive rotations by angles
$\alpha_j\in(0,\pi)$.  In one coordinate use a half-disk with negative
puncture at zero and the two boundary conditions on its real radii.
Writing $a=\alpha_j/\pi$, the section $z^a$ decays at the puncture.
Choose the plane path on the rounded outer boundary with angle
$a\arg z$.
Dividing by $z^a$ removes the rotating boundary condition and leaves the
real disk operator with a removable boundary point and a free real
constant mode.  Its index is one.  Equivalently, moving the end weight
to the removable-point problem crosses no other indicial root.  Thus
the minimal positive rotations contribute eight; further winding adds
its Maslov index.  With the grading lifts $\vartheta_1,\vartheta_2$ this
gives
\[
 8+\vartheta_2-\vartheta_1-\sum_j\alpha_j/\pi.
\]
The reciprocal output cap has complementary index, so the two cap
indices sum to eight and have the same parity.

After capping all transverse ends, the disk Maslov index is a multiple
of four; it is zero for the integral matching on a graph word.
Fredholm indices add under transverse capping.  Adding the $k-2$
domain parameters for $k+1$ boundary marks and then taking the
evaluation pushforward gives degree $2-k$ with the indicated grading.
For instance, with jumps only, if the input degrees are $d_j$ and
the output degree is $d_0$, the caps have total index
$\sum_jd_j+8-d_0$.  The uncapped operator therefore has index
$\mu+d_0-\sum_jd_j$; virtual dimension zero gives
$d_0=\sum_jd_j+2-k-\mu$.  Inserting a smooth form and taking an
evaluation pushforward gives the same formula with its form degree
in place of the jump degree.
For a pure disk the uncapped real disk index is $8+\mu$.
There are $k+1$ real marked-point parameters in the boundary-output
case and $k+2$ in the interior-output case, and the disk automorphism
group has dimension three.  This proves both dimension formulas.
In the pure case the ordinary form degree of the output is
\[
 \sum_{j=1}^k|x_j|+2-k-\mu.
\]
Putting every input in degree $1$ modulo four gives degree $2$ modulo
four.  On the eight-manifold $L$ these are exactly degrees $2$ and $6$.
\end{proof}

Here is the determinant order used for the geometric correspondences.
At a smooth slot put $p_j=0$, and at a jump put $p_j$ equal to the
generator degree modulo two, represented by $0$ or $1$.  Let $X_j$ be
the degree of the form at that slot, zero at jumps; $X_0$ denotes the
degree of an output test form when there is one.  Put
\begin{equation}
 \label{obs:determinant-order}
 K W_0W_1\cdots W_k,\qquad
 K=o(D_{\mathrm{cap}})\,o(\mathrm{PSL}(2,\R))^*,\qquad
 W_j=U_jC_j^*.
\end{equation}
Here $U_j$ is the oriented tangent of the mark, positive
counterclockwise at a boundary mark, and $C_j$ is the cap determinant
factor, trivial at a smooth slot.  Thus $K$ has odd parity and a
boundary factor $W_j$ has parity $1+p_j$.  For an interior-output
operation $W_0$ is just the even tangent factor of the interior mark.
All rearrangements of determinant lines use the Koszul rule.

Consider a normalized boundary face that splits off $s$ consecutive
inputs, beginning at slot $i$.  Let $D'$ be the child and $A'$ the
parent.  Their respective arities are $s$ and
$b=k-s+1$.  Let $z$ denote the new node slot, and put
\[
 l=i-1,\qquad r=k-i-s+1.
\]
Subscripts $<$, $D'$, and $>$ on $P$ or $X$ denote sums over the
preceding, child, and trailing input blocks; $P$ is always a sum of
the $p_j$.  The case $s=0$ is included, with an empty child input block.

\begin{lemma}[Boundary orientation]\label{obs:boundary-sign}
Give a product face the order parent, child, followed at a smooth node
by the dual tangent of the matching copy of $L$.  The boundary
orientation relative to this product has exponent
\begin{equation}
 \label{obs:boundary-exponent}
 \sigma=i+p_0+P_<+
       (p_z+s+P_{D'})(r+P_>)+\varepsilon+\rho_z.
\end{equation}
Here $\varepsilon$ is one fixed smooth disk gluing sign, and $\rho_z$
is a fixed additional contraction sign for the chosen transverse cap
pair at $z$; it is zero at a smooth node.  For an interior-output
parent and a nonwrapping placement of a pure child, the exponent is
\begin{equation}
 \label{obs:p-boundary-exponent}
 i-1+sr+\varepsilon.
\end{equation}
The same signs hold for the thickened charts and their normalized
faces.
\end{lemma}

\begin{proof}
We first justify that the residual contraction is local.  Cap the
exterior jumps before gluing the distinguished node, using
associativity of determinant gluing.  At a smooth node the spin
boundary framings concatenate and one contracts against the oriented
tangent diagonal.  At a transverse node one compares strip gluing
with gluing the two reciprocal caps at that node.  The extra
determinant pairing therefore uses only this fixed cap pair.

More explicitly, choose the spin framing on a boundary collar and
fill that collar by a trivial real disk.  The remaining clutching is
transferred to a complex sphere away from the distinguished end.
This is the excision construction of the spin disk orientation.
The sphere determinant and its interior matching factor are
complex-oriented even lines, so moving them past the node
contraction creates no sign.  The boundary framing after gluing is
the concatenation in its spin homotopy class, with the node cap
reversed when appropriate.  Thus changes in exterior boundary
paths or in their Maslov contributions do not alter this local
comparison.  Its signs are the fixed $\varepsilon$ and $\rho_z$
in the statement.

The domain factor is local for the same reason.  Put the child node
at infinity in an upper half-plane coordinate and insert the child
at a parent mark $a$ by $\zeta=a+\rho z$, $\rho>0$.
Use the outward normal, directed toward the broken configuration,
first.  The child's residual affine automorphisms are translation
and dilation.  They cancel with the parent-node position and the
radial smoothing parameter.  In a cross-section to these two
directions all retained real mark tangents are transported by
positive scaling, while interior-mark pairs are transported with
their complex orientations.  One can make this calculation before
quotienting the parent automorphisms and then quotient by the same
parent directions on both sides.  The sign is consequently
independent of the other marks, including the presence of an
interior mark on the parent.  Changing upper half-plane charts
preserves these orientations.  Auxiliary interior stabilization
marks and slice constraints add matching even-dimensional factors
and do not change the comparison.

It remains to perform the permutation in
Equation~\eqref{obs:determinant-order}.  Write
\[
 \mathsf A=1+p_0+l+P_<,\qquad
 \mathsf R=r+P_>,\qquad \mathsf C=s+P_{D'}.
\]
Starting with outward normal, parent factors, child factors, and
the smooth matching factor when present, group the local factors as
\[
 N'K_{A'}W_{A'z}K_{D'}W_{D'z}.
\]
The parent node factor crosses the preceding parent slots, the
child $K$ crosses the remaining parent slots, and the child node
factor crosses those same slots.  Finally the child input block
crosses the trailing parent slots.  Their total permutation sign is
\begin{align*}
 &(1+p_z)\mathsf A+(\mathsf A+\mathsf R)
       +(1+p_z)(\mathsf A+\mathsf R)+\mathsf C\mathsf R\\
 &\hspace{20mm}=\mathsf A+(p_z+\mathsf C)\mathsf R
   =i+p_0+P_<+(p_z+s+P_{D'})(r+P_>).
\end{align*}
The matching tangent has dimension eight and contributes no further
permutation parity.  This proves
Equation~\eqref{obs:boundary-exponent}.  For an interior-output
parent all boundary slots are smooth and its output factor is
even, so the same computation has $\mathsf A=l$, $\mathsf R=r$,
$\mathsf C=s$, and
$p_z=0$.  It gives Equation~\eqref{obs:p-boundary-exponent}.

For the thickening, append the oriented obstruction coefficient
directions to the Fredholm determinant.  They have even dimension
and split by components on the face, so the Fredholm and tangent
exact sequences give the same permutation.  Compact lower-order
terms in the linearized equations, including derivatives of the
map-dependent core insertions, can be homotoped to zero with fixed
asymptotic operators and compatibly with the split face.
Determinant orientations persist under this homotopy.  True chart
isomorphisms preserve the cappings, spin data, and external slots;
stabilization pairs and permutations of probe blocks are even.
Their orientation compatibility on the smooth locus therefore
extends across the corners by this same gluing comparison, also
when components and their smoothing parameters are permuted.
\end{proof}

\subsection{Normalization of the operations}

For a boundary word $W$ of arity $k$, put
\begin{equation}
 \label{obs:normalization}
 \begin{split}
 e_W&=k+P-1-p_0,\\
 G_W&=\binom{e_W}{2}+p_0e_W+P+\varepsilon
                  +\sum_{\text{input jumps }j}\rho_j,\\
 S_W&=\sum_{j=1}^k(j+P_{<j})X_j+G_W.
 \end{split}
\end{equation}
The expression $\binom e2=e(e-1)/2$ is used for every integer $e$,
including negative integers.  This matters for zero-input terms.
Define a stable $m$ term by multiplying the geometric integral by
$(-1)^{S_W}$, with the pulled-back forms ordered as output test
first, then inputs.  Thus at a smooth output its definition is
\[
 \int_L y_0\wedge m_{k,\beta}(x_1,\ldots,x_k)
 =(-1)^{S_W}
   \int_{\mathcal M_{k,\beta}}^{\mathrm{vir}}
          \mathrm{ev}_0^*y_0\wedge
          \mathrm{ev}_1^*x_1\wedge\cdots\wedge
          \mathrm{ev}_k^*x_k,
\]
and at a jump output it defines the scalar coefficient of that
generator, omitting the output form.  We sum these terms with their
energy monomial and, on the permitted words, the relative weight
$\exp(tB_{\mathrm{rel}})$.  Add the unstable pure unary operation
$m_{1,0}=d$.

For a pure interior-output operation tested by $\eta\in A^*(T)$,
write $p_k^\eta(x_1,\ldots,x_k)$ for the resulting scalar integral.
Its operation sign is
\begin{equation}
 \label{obs:p-normalization}
 \sum_{j=1}^k jX_j+\binom{k-1}{2}+\varepsilon,
\end{equation}
again with the output test form first.  There is no unstable
constant $p_0$ term.  The stable constant $p_{1,0}$ term is retained.

\begin{proposition}[Finite signed identities]\label{obs:ainfinity}
Within the finite cost system the normalized operations satisfy
the curved suspended identities
\begin{equation}
 \label{obs:ainfinity-equation}
 \sum_{s=0}^k\ \sum_{i=1}^{k-s+1}
 (-1)^{\sum_{j<i}(|x_j|-1)}
 m_{k-s+1}\bigl(x_1,\ldots,x_{i-1},
       m_s(x_i,\ldots,x_{i+s-1}),
       x_{i+s},\ldots,x_k\bigr)=0.
\end{equation}
For a closed interior test form, the corresponding Stokes identity
for $p^\eta$ has the same insertion signs and no output
differential.  These assertions hold coefficientwise in $t$ with
the relative weight, as well as without that weight.
\end{proposition}

\begin{proof}
The virtual Stokes and product formulas are exact by
Theorem~\ref{virt:finite-integration}.  We verify their normalization.
The parity condition for an $m$ Stokes integral is
\[
 X_0+\sum_jX_j=p_0+P+k+1.
\]
Multiply its Stokes equation by $(-1)^{S_W+X_0+1}$.
Integration by parts turns the $dy_0$ term into the output
differential with coefficient one, because
\[
 \int_Ldy_0\wedge u=(-1)^{X_0+1}\int_Ly_0\wedge du.
\]
For an input differential $dx_i$, the change of $S_W$ is
$i+P_{<i}$, while the Leibniz sign before that input is
$X_0+X_{<i}$.  The normalized exponent is therefore
\[
 1+X_{<i}+i+P_{<i}
 =\sum_{j<i}(p_j+X_j-1)
 =\sum_{j<i}(|x_j|-1).
\]

Consider next the split in Lemma~\ref{obs:boundary-sign}.  Let
$X_z$ be the degree of the \emph{output} form after the child
pushforward, not the degree of a test form on that output.  The
degree formula gives
\[
 X_{D'}+X_z=s+P_{D'}+p_z=:u.
\]
The Fubini comparison contributes $uX_>$: move the child input
forms past the trailing forms and then move its output form back
into the node slot.  Integration uses the test form first on the
even-dimensional matching base.  For the form-dependent parts
$\Phi_W=S_W-G_W$, the preceding input coefficients cancel,
the child input coefficients differ by $l+P_<$, and the
trailing coefficients differ by $u+1$.  Hence
\[
 \Phi_W+\Phi_{A'}+\Phi_{D'}
 =(l+P_<)X_{D'}+(l+1+P_<)X_z+(u+1)X_>.
\]
After adding $uX_>$ this becomes
$(l+P_<)u+X_z+X_>$.  Use
$X_z=X_{D'}+u$ and the Stokes parity equation above.
Subtracting the boundary term
from Stokes and including $S_W,S_{A'},S_{D'}$ and
Equation~\eqref{obs:boundary-exponent}, the difference from the
desired insertion exponent is
\begin{equation}
 \label{obs:sign-residual}
 k+P+u(b+P_<+P_>)+\varepsilon+\rho_z
          +G_W+G_{A'}+G_{D'}.
\end{equation}
We check this cancellation without specifying the bare signs.
Set
\[
 a_0=p_0,\quad q=p_z,\quad
 v=l+P_<,\quad w=r+P_>,\quad c=s+P_{D'}.
\]
Then
\[
 e_{A'}=v+w+q-a_0,\qquad e_{D'}=c-1-q,
 \qquad e_W=e_{A'}+e_{D'}.
\]
The binomial terms in the sum of the three $G$'s give
$e_{A'}e_{D'}$.  The terms involving the output parity give
$(a_0+q)e_{D'}$, the three input-parity sums give $q$, and the
remaining constants give $\varepsilon+\rho_z$: every exterior
input contraction sign occurs twice, while the node sign occurs
once.  Consequently Equation~\eqref{obs:sign-residual} reduces to
\[
 v+w+c+(c+q)(v+w+1)+(v+w)(c-1-q)+q=0
 \pmod2.
\]
This verifies the split term, including $s=0$ and negative values
of $e_{D'}$.

For pure $p$ the parity condition for the Stokes integral changes
by one, to $X_0+\sum X_j=k$, and the boundary exponent changes by
one, from $i+sr+\varepsilon$ to
$i-1+sr+\varepsilon$.  These changes cancel in the preceding
normalization.  The input-differential computation is unchanged.
There is no $d\eta$ contribution for a closed test form.  This
proves the asserted interior-output identity on nonwrapping
blocks; cyclic reindexing supplies the other blocks, with the
sign computed below.

Finally, Lemma~\ref{local:weight} gives constancy of
$B_{\mathrm{rel}}$ on each retained relative class and its
additivity on a split.  Thus
$\exp(tB_{\mathrm{rel}})$ is constant for Stokes and is the
product of the two component weights for Fubini.  It introduces
no extra differential or boundary term.  The same proof therefore
works separately at every Taylor coefficient.
\end{proof}

\begin{lemma}[Cyclic sign]\label{obs:cyclic}
Moving the first boundary input of $p_k^\eta$ to the last slot
multiplies the operation by
\begin{equation}
 \label{obs:cyclic-sign}
 (-1)^{(X_1-1)\sum_{j>1}(X_j-1)}.
\end{equation}
This is a reindexing identity on the labeled moduli spaces and
does not impose a further quotient by cyclic rotations.
\end{lemma}

\begin{proof}
The boundary tangent factors give $k-1$, permuting the input
forms gives $X_1\sum_{j>1}X_j$, and changing
Equation~\eqref{obs:p-normalization} gives
$(k-1)X_1+\sum_{j>1}X_j$.  Their sum is
$(X_1-1)\sum_{j>1}(X_j-1)$ modulo two.  The interior mark is even
and stays first.  Cyclic equivariance of the raw integration is
part of Theorem~\ref{virt:finite-integration}; hence this
calculation also applies to its normalized faces and finite
isotropy factors.
\end{proof}

\subsection{The interior-output identity}

Let $b$ be a formal cochain with coefficients in
$A^{\bar1}(L)[[t]]$, only finitely many energy exponents at
most $N$, and all energies at least $\delta_L$, the pure
disk gap from Lemma~\ref{local:gap}.  Set
\begin{equation}
 \label{obs:curvature}
 \Delta(b)=\sum_{k\ge0}m_k(b,\ldots,b),
\end{equation}
including $db$ in the unary term.  Through energy $N$ the
sum is finite.  Every $b$ input is even after suspension.

\begin{proposition}[Obstruction pairing identity]
\label{obs:pairing-identity}
For every closed form $\eta$ on $T$, the finite signed
correspondences give
\begin{equation}
 \label{obs:weighted-stokes}
 \sum_{r\ge0}p_{r+1}^\eta(\Delta(b),b,\ldots,b)\equiv_N0.
\end{equation}
The statement uses the weighted operations on the pure $L$
sector and holds coefficientwise in $t$.
\end{proposition}

\begin{proof}
For $k>0$ apply the interior-output Stokes identity to $k$
copies of $b$ and divide by $k$.  A child carrying a block of
$s$ of the $k=r+s$ inputs occurs in $k$ cyclic placements.
For $0<s<k$ these are the choices of the start of the block.
For $s=0$ the zero-input bubble attaches in any of the $k$
cyclic gaps.  For $s=k$ all inputs lie on the child, but its
output node can still occupy any of those $k$ gaps; these give
the $k$ possible linear input orders on the child.  Thus the
factor $k$ occurs also at both extremes.  An input differential
likewise occurs in $k$ slots.  After rotating its resulting
curvature slot to the first position, all signs are positive:
the other inputs have suspended parity zero, so
Equation~\eqref{obs:cyclic-sign} contributes no sign, even
though the curvature itself has suspended parity one.

For $k=0$ use coefficient one.  There is one normalized
boundary type $p_1^\eta(m_0)$, with the interior mark
distinguishing the parent.  Its constant parent is stable,
having one interior mark and one boundary node.  This argument
does not add the nonexistent unstable constant $p_0$.
The usual finite-isotropy integration already counts a face
with its distinguished node; there is no extra rotation
factor.  In particular a finite-order disk automorphism
fixing that boundary node is the identity on the incident
disk, and the distinction propagates along attached disks.
If an automorphism of a $p_0$ tree permutes indistinguishable
edges before the edge is distinguished, its ordinary orbifold
factor and the normalized-boundary multiplicity give exactly
the preceding count.  Summing these equations over $k$ gives
Equation~\eqref{obs:weighted-stokes}.

All equations used lie within the finite cost cut.  At total
energy at most $N$ there are at most $N/\delta_L$ inserted
$b$'s.  Their arity cost is less than $N$, since
$a<\delta_L$ in Proposition~\ref{local:finite}.  The
interior-output equation therefore has cost less than $2N$,
and splitting preserves cost.  This is contained in the
chosen cut of $4N$.
\end{proof}

\begin{theorem}[Finite weighted bounding cochains]
\label{obs:finite-bounding}
For every integer $N\ge1$ and every finite list of graph labels, the
finite operations can be chosen with a cochain
\[
 b_{L,N}(t)=\sum_{\lambda\in F_N}\qnov^\lambda b_\lambda(t),
 \qquad F_N\subset[\delta_L,N]\ \text{finite},\qquad
 b_\lambda(t)\in A^{\bar1}(L)[[t]],
\]
such that $\Delta(b_{L,N})\equiv_N0$.
This assertion holds for all
Taylor coefficients simultaneously.  Only a finite order of
differentiability, chosen sufficiently high depending on $N$
and the finite list, is needed.  No compatibility as $N$
varies is asserted or required.
\end{theorem}

\begin{proof}
Begin with $b=0$ and proceed in increasing positive energy.
Suppose all curvature coefficients below $\lambda$ have
been canceled.  Summing
Equation~\eqref{obs:ainfinity-equation} on strings of $b$'s
gives the Bianchi identity
\[
 \sum_{u,v\ge0}m_{u+v+1}(b^{\otimes u},\Delta(b),
                              b^{\otimes v})\equiv_N0.
\]
At energy $\lambda$, every nonconstant operation or additional
$b$ input raises the energy.  The only term involving the
first uncanceled curvature is therefore the constant unary
term.  It follows that
\begin{equation}
 \label{obs:closed-obstruction}
 d\Delta_\lambda(t)=0.
\end{equation}
The identical energy argument applied to
Equation~\eqref{obs:weighted-stokes} leaves only
$p_{1,0}^\eta(\Delta_\lambda(t))$.  Its moduli space is the
ordinary copy of $L$, with its boundary evaluation the
identity and its interior evaluation $i_L$.
The constant disk has $B_{\mathrm{rel}}=0$.
By Theorem~\ref{virt:finite-integration} its raw integral is
ordinary integration.  For the even-degree form
$\Delta_\lambda$ the sign in
Equation~\eqref{obs:p-normalization} is one fixed sign,
independent of $\lambda$, of its Taylor coefficient, and of
the choice of complementary even-degree test form.  Hence
\begin{equation}
 \label{obs:ambient-annihilation}
 \int_Li_L^*\eta\wedge\Delta_\lambda(t)=0.
\end{equation}
Only its nonzero sign, not a numerical convention for it, is
needed for this implication.

By Lemma~\ref{obs:degree}, write the curvature coefficient as
the sum of its degree-two and degree-six parts.
For a closed degree-two coefficient $u$, Equation~\eqref{obs:ambient-annihilation}
for all closed ambient six-forms says that
$i_{L*}(u\cap[L])\in H_6(T;\C)$ pairs to zero with all of
$H^6(T;\C)$, and is therefore zero.  The injection on $H_6$
in Equation~\eqref{obs:topological-input} makes $u\cap[L]=0$.
Poincar\'e duality on
$L$ gives $[u]=0$.  For a degree-six coefficient the same
argument uses closed ambient two-forms and injectivity on
$H_2$.  Rational injectivity remains injectivity after scalar
extension to $\C$.  Applying this argument to each Taylor
coefficient proves that $\Delta_\lambda(t)$ is exact
coefficientwise.

Fix a Riemannian metric on $L$ and a linear primitive
operator on exact forms, for example $h=d^*G$ with $G$ the
Hodge Green operator.  For a closed exact form $u$ it
satisfies $dhu=u$.  Replace
\[
 b\quad\text{by}\quad b-\qnov^\lambda h\Delta_\lambda(t).
\]
The correction has degrees one and five.  Its unary
differential cancels precisely the energy-$\lambda$
curvature.  Every other new contribution has greater
energy: a nonconstant vertex has positive area, and a
constant stable $m$ vertex has at least two inputs, so any
such occurrence of the correction is accompanied by
another positive-energy input.  Thus no earlier coefficient
is changed.  The lower energy bound $\delta_L$ is preserved.

For completeness, this recursion is finite without assuming
that the symplectic period group is discrete.  There are
only finitely many positive energies of nonconstant pure
terms in the relevant cut by Proposition~\ref{local:finite};
call their set $\mathcal E_N$.  Each is at least
$\delta_L$.  Every energy introduced by a curvature
operation or by taking its primitive is a sum of energies
in $\mathcal E_N$.  Below $N$ such a sum has at most
$\lfloor N/\delta_L\rfloor$ summands, so only finitely many
sums occur.  Zero-area stable vertices add no new
exponents.  In the expression trees, after ignoring the
linear $d$ and $h$ edges, every such vertex has at least two
inputs; their number is bounded by the number of
positive-energy vertices and the external arity.  Thus
only finitely many primitive and composition steps are
needed.

This bound is independent of Taylor order.  The formal
weight changes the scalar at a vertex; coefficientwise
products in a fixed finite arity use finite Cauchy sums.
The same finite expression trees and the same primitive
operator handle every Taylor coefficient.  Choose the
finite chart regularity in Proposition~\ref{local:charts}
high enough for the bounded number of pushforwards,
compositions, differentials, and primitives.  Their finite
losses of regularity then leave the regularity needed for
the identities throughout the recursion.  This proves the
simultaneous formal assertion and completes the induction
through energy $N$.
\end{proof}

\subsection{Twisted mixed operations}

Write $b=b_{L,N}(t)$.  For a strictly directed graph word
$X_0\prec\cdots\prec X_d$, graph inputs $x_1,\ldots,x_d$,
and $y\in\C\langle X_d\cap L\rangle$, define
\begin{equation}
 \label{obs:module-twist}
 \mu^{b}_{d\mid1}(x_1,\ldots,x_d;y)
 =\sum_{\ell\ge0}m_{d+1+\ell}
       (x_1,\ldots,x_d,y,b^{\otimes\ell}),
\end{equation}
truncated to total energy at most $N$.  Its degree is $1-d$
modulo four.  The next result shows that $\mu^b_{0\mid1}$
is a differential through the cutoff on the finite free
intersection-generator space at $X_0$.

\begin{corollary}[Finite module identities]
\label{obs:finite-module}
For external graph arity at most $N$, the operations in
Equation~\eqref{obs:module-twist} satisfy the right module
identities through energy $N$ over the directed graph
operations.  All coefficient entries have nonnegative
energy, and the errors in these finite matrix identities
have energy greater than $N$, coefficientwise in $t$.
On words without a $C$-slope label the operations use the
relative weight; their specialization at $t=0$ is exactly
the restriction of the unweighted operations on the full
finite list, including its $C$-slope labels, twisted by
$b_{L,N}(0)$.
\end{corollary}

\begin{proof}
Apply Equation~\eqref{obs:ainfinity-equation} to the mixed
words with the appended $b$'s and sum over their number.
The $b$'s have suspended parity zero, so these summations
introduce no new insertion signs.  A split entirely in the
graph inputs gives a directed category operation.  A mixed
child gives a composition of two operations $\mu^b$.
The remaining pure splits, together with the differentials
of appended $b$'s, sum to an insertion of $\Delta(b)$ and
vanish through energy $N$.  Explicitly, putting
$\gamma_j=|x_j|-1$, the remaining identity is
\begin{align*}
0\equiv_N{}&
 \sum_{s=1}^{d}\ \sum_{i=1}^{d-s+1}
 (-1)^{\sum_{j<i}\gamma_j}
 \mu^b_{d-s+1\mid1}
    (x_1,\ldots,m_s(x_i,\ldots,x_{i+s-1}),\ldots,x_d;y)\\
 &+\sum_{i=0}^{d}(-1)^{\sum_{j\le i}\gamma_j}
 \mu^b_{i\mid1}
    (x_1,\ldots,x_i;
       \mu^b_{d-i\mid1}(x_{i+1},\ldots,x_d;y)).
\end{align*}
The case $d=0$ states $(\mu^b_{0\mid1})^2\equiv_N0$.
There is no directed curvature term, because the graphs
are weakly exact and all their nonidentity words are
strictly directed.

At energy at most $N$ the twist uses at most
$N/\delta_L$ appended copies of $b$.  If $d\le N$, its
stable $m$ cost is bounded by
\[
 N+a(d+\ell)\le N+a(N+N/\delta_L)<3N.
\]
The same bound applies to the total cost of a split by
additivity.  Every term of the asserted identity is
therefore present in the $4N$ system.  Discarding terms
above $N$ cannot change any coefficient at or below $N$,
as operations and $b$ have nonnegative energies.  On the
jump sector the generator sets are finite, so these are
finite matrix identities in each fixed external arity.

The weighted and unweighted constructions use the same
finite geometric integrations and signs.  Pure $L$ terms
admit the relative weight, and setting $t=0$ in the
constructed bounding cochain cancels their unweighted
curvature.  The unweighted mixed identities then apply
also to words containing $C$-slope labels.  On the common
words the weight specializes to one, proving the claimed
compatibility.
\end{proof}
